\textbf{Phần 1. Trắc nghiệm lựa chọn}
\setcounterpageref{ex}{0}

\begin{ex} %[5D4N24]
Kết quả của phép tính $52 - 4{,}91$ là
\choice
{\True $47{,}09$}
{$47{,}19$}
{$48{,}09$}
{$47{,}91$}
\loigiai{
Ta có: $52{,}00 - 4{,}91 = 47{,}09$.\\
\textbf{Đáp án của câu hỏi là A}
}
\end{ex}

\begin{ex} %[5D4N24]
Kết quả của phép tính $0{,}73 \times 6{,}5$ là
\choice
{$47{,}45$}
{$0{,}4745$}
{$4{,}755$}
{\True $4{,}745$}
\loigiai{
Ta có: $0{,}73 \times 6{,}5 = 4{,}745$.\\
\textbf{Đáp án của câu hỏi là D}
}
\end{ex}

\begin{ex} %[5D4H24]
Kết quả của phép tính $104{,}89 : 8{,}5$ là
\choice
{$123{,}4$}
{$1{,}234$}
{\True $12{,}34$}
{$12{,}44$}
\loigiai{
Ta có: $104{,}89 : 8{,}5 = 1\,048{,}9 : 85 = 12{,}34$.\\
\textbf{Đáp án của câu hỏi là C}
}
\end{ex}

\begin{ex} %[5D4H24]
Giá trị của biểu thức $(131{,}4 - 80{,}8) : 2{,}3 + 21{,}64 \times 2$ là
\choice
{\True $65{,}28$}
{$44{,}28$}
{$55{,}28$}
{$75{,}28$}
\loigiai{
Ta có: $50{,}6 : 2{,}3 + 43{,}28 = 22 + 43{,}28 = 65{,}28$.\\
\textbf{Đáp án của câu hỏi là A}
}
\end{ex}

\begin{ex} %[5D4H24]
Tính thuận tiện biểu thức $0{,}56 \times 4{,}7 + 5{,}3 \times 0{,}56$ ta được kết quả là
\choice
{$0{,}56$}
{$56$}
{$560$}
{\True $5{,}6$}
\loigiai{
Ta có: $0{,}56 \times (4{,}7 + 5{,}3) = 0{,}56 \times 10 = 5{,}6$.\\
\textbf{Đáp án của câu hỏi là D}
}
\end{ex}

\begin{ex} %[5D4H24]
Tìm $x$ biết: $51{,}23 + x = 74{,}9$.
\choice
{$23{,}77$}
{$22{,}67$}
{\True $23{,}67$}
{$126{,}13$}
\loigiai{
Ta có: $x = 74{,}90 - 51{,}23 = 23{,}67$.\\
\textbf{Đáp án của câu hỏi là C}
}
\end{ex}

\begin{ex} %[5D4H24]
Lượng mưa đo được trong $3$ lần của tháng Hai lần lượt là $15{,}6\text{ mm}$, $18{,}9\text{ mm}$ và $12{,}3\text{ mm}$. Lượng mưa trung bình của tháng Hai là
\choice
{$14{,}6\text{ mm}$}
{\True $15{,}6\text{ mm}$}
{$15{,}3\text{ mm}$}
{$16{,}2\text{ mm}$}
\loigiai{
Lượng mưa trung bình là: $(15{,}6 + 18{,}9 + 12{,}3) : 3 = 46{,}8 : 3 = 15{,}6\text{ (mm)}$.\\
\textbf{Đáp án của câu hỏi là B}
}
\end{ex}

\begin{ex} %[5D4V24]
Dúi làm rơi khúc mía xuống một cái lỗ sâu $1{,}31\text{ m}$ ($1{,}06\text{ m} + 0{,}25\text{ m}$). Chồn đã đào một cái hố sâu $0{,}9\text{ m}$ bên cạnh để lấy giúp. Hỏi chồn còn phải đào sâu thêm bao nhiêu mét nữa?
\choice
{$0{,}31\text{ m}$}
{$0{,}51\text{ m}$}
{\True $0{,}41\text{ m}$}
{$2{,}21\text{ m}$}
\loigiai{
Độ sâu cần đào thêm là: $1{,}31 - 0{,}9 = 0{,}41\text{ (m)}$.\\
\textbf{Đáp án của câu hỏi là C}
}
\end{ex}

\begin{ex} %[5D4V24]
Một tấm bản đồ hình chữ nhật có chiều dài $8{,}4\text{ dm}$ và chiều rộng $6{,}5\text{ dm}$. Sau khi gấp đôi $2$ lần (thành $4$ phần bằng nhau) ta được một hình chữ nhật nhỏ có diện tích là
\choice
{$54{,}6\text{ dm}^2$}
{\True $13{,}65\text{ dm}^2$}
{$27{,}3\text{ dm}^2$}
{$13{,}5\text{ dm}^2$}
\loigiai{
Diện tích bản đồ là: $8{,}4 \times 6{,}5 = 54{,}6\text{ (dm}^2\text{)}$.\\
Diện tích hình chữ nhật nhỏ sau $2$ lần gấp đôi là: $54{,}6 : 4 = 13{,}65\text{ (dm}^2\text{)}$.\\
\textbf{Đáp án của câu hỏi là B}
}
\end{ex}

\begin{ex} %[5D4V24]
Trong hai năm, trang trại bán được tất cả $21{,}56\text{ tấn}$ cá chẽm. Biết năm thứ hai bán được nhiều hơn năm thứ nhất $2{,}7\text{ tấn}$. Số tấn cá chẽm trang trại bán được trong năm thứ nhất là
\choice
{\True $9{,}43\text{ tấn}$}
{$12{,}13\text{ tấn}$}
{$10{,}78\text{ tấn}$}
{$8{,}43\text{ tấn}$}
\loigiai{
Năm thứ nhất bán được: $(21{,}56 - 2{,}7) : 2 = 18{,}86 : 2 = 9{,}43\text{ (tấn)}$.\\
\textbf{Đáp án của câu hỏi là A}
}
\end{ex}

\begin{ex} %[5D4V24]
Khu vườn trồng cây ăn quả dạng hình chữ nhật có diện tích $83{,}52\text{ m}^2$ và chiều rộng $8{,}7\text{ m}$. Chu vi khu vườn đó là
\choice
{$18{,}3\text{ m}$}
{\True $36{,}6\text{ m}$}
{$34{,}8\text{ m}$}
{$73{,}2\text{ m}$}
\loigiai{
Chiều dài khu vườn là: $83{,}52 : 8,7 = 9{,}6\text{ (m)}$.\\
Chu vi khu vườn là: $(9{,}6 + 8{,}7) \times 2 = 18{,}3 \times 2 = 36{,}6\text{ (m)}$.\\
\textbf{Đáp án của câu hỏi là B}
}
\end{ex}

\begin{ex} %[5D4V24]
Bạn Tú có $25$ quả bóng, mỗi quả nặng $5{,}7\text{ g}$. Tú bỏ tất cả số bóng vào hai hộp xanh và đỏ. Tú thấy cân nặng của hộp xanh tăng thêm $62{,}7\text{ g}$. Hỏi cân nặng của hộp đỏ tăng thêm bao nhiêu gam?
\choice
{$142{,}5\text{ g}$}
{$74{,}1\text{ g}$}
{$80{,}2\text{ g}$}
{\True $79{,}8\text{ g}$}
\loigiai{
Tổng khối lượng $25$ quả bóng là: $25 \times 5{,}7 = 142{,}5\text{ (g)}$.\\
Cân nặng hộp đỏ tăng thêm là: $142{,}5 - 62{,}7 = 79{,}8\text{ (g)}$.\\
\textbf{Đáp án của câu hỏi là D}
}
\end{ex}

\textbf{Phần 2. Trắc nghiệm đúng sai}
\setcounterpageref{ex}{0}

\begin{ex} %[5D4H24]
Xét tính đúng sai của các khẳng định sau:
\choiceTF
{\True Phép tính $4{,}98 + 15{,}7 = 20{,}68$.}
{Phép tính $314 : 0{,}001 = 0{,}314$.}
{\True Biểu thức $(5{,}3 + 4{,}7) \times 0{,}1 = 5{,}3 \times 0{,}1 + 4{,}7 \times 0{,}1 = 1$.}
{\True Giá trị biểu thức $8{,}16 : (1{,}32 + 3{,}48) - 0{,}34 : 2 = 1{,}53$.}
\loigiai{
\begin{itemize}
    \item $4{,}98 + 15{,}70 = 20{,}68$.
    \item $314 : 0{,}001 = 314 \times 1\,000 = 314\,000 \neq 0{,}314$.
    \item $(5{,}3 + 4{,}7) \times 0{,}1 = 10 \times 0{,}1 = 1$.
    \item $8{,}16 : 4{,}8 - 0{,}17 = 1{,}7 - 0{,}17 = 1{,}53$.
\end{itemize}
\begin{itemize}
    \item a) \textbf{Đúng}.
    \item b) \textbf{Sai}: Kết quả là $314\,000$.
    \item c) \textbf{Đúng}.
    \item d) \textbf{Đúng}.
\end{itemize}
}
\end{ex}

\begin{ex} %[5D4V24]
Một mảnh đất hình chữ nhật có kích thước $40{,}4\text{ m} \times 30\text{ m}$. Người ta sử dụng đúng một nửa diện tích mảnh đất để xây nhà máy. Xét tính đúng sai của các khẳng định sau:
\choiceTF
{\True Diện tích của cả mảnh đất là $1\,212\text{ m}^2$.}
{Diện tích đất dùng để xây nhà máy là $600\text{ m}^2$.}
{\True Diện tích nhà máy bằng $0{,}0606\text{ ha}$.}
{\True Diện tích của nhà máy lớn hơn $\dfrac{1}{100}\text{ ha}$ (tức $100\text{ m}^2$).}
\loigiai{
\begin{itemize}
    \item Diện tích mảnh đất: $40{,}4 \times 30 = 1\,212\text{ m}^2$.
    \item Diện tích xây nhà máy: $1\,212 : 2 = 606\text{ m}^2 \neq 600\text{ m}^2$.
    \item $606\text{ m}^2 = 606 : 10\,000 = 0{,}0606\text{ ha}$.
    \item $\dfrac{1}{100}\text{ ha} = 100\text{ m}^2 < 606\text{ m}^2$.
\end{itemize}
\begin{itemize}
    \item a) \textbf{Đúng}.
    \item b) \textbf{Sai}: Diện tích là $606\text{ m}^2$.
    \item c) \textbf{Đúng}.
    \item d) \textbf{Đúng}.
\end{itemize}
}
\end{ex}

\textbf{Phần 3. Trả lời ngắn}
\setcounterpageref{ex}{0}

\begin{ex} %[5D4H24]
Tính giá trị biểu thức: $(131{,}4 - 80{,}8) : 2{,}3 + 21{,}64 \times 2$. (Nhập kết quả dưới dạng số thập phân).
\shortans{65,28}
\loigiai{
\textbf{Đáp số:} $65{,}28$\\
$(131{,}4 - 80{,}8) : 2{,}3 + 21{,}64 \times 2 = 50{,}6 : 2{,}3 + 43{,}28 = 22 + 43{,}28 = 65{,}28$.
}
\end{ex}

\begin{ex} %[5D4H24]
Tính nhanh giá trị biểu thức: $0{,}56 \times 4{,}7 + 5{,}3 \times 0{,}56$. (Nhập kết quả dưới dạng số thập phân).
\shortans{5,6}
\loigiai{
\textbf{Đáp số:} $5{,}6$\\
$0{,}56 \times (4{,}7 + 5{,}3) = 0{,}56 \times 10 = 5{,}6$.
}
\end{ex}

\begin{ex} %[5D4H24]
Tìm $x$ biết: $x : 3{,}8 = 21{,}34$. (Nhập kết quả dưới dạng số thập phân).
\shortans{81,092}
\loigiai{
\textbf{Đáp số:} $81{,}092$\\
$x = 21{,}34 \times 3{,}8 = 81{,}092$.
}
\end{ex}

\begin{ex} %[5D4V24]
Tấm bản đồ hình chữ nhật có kích thước $8{,}4\text{ dm} \times 6{,}5\text{ dm}$. Tính diện tích tấm bản đồ đó theo đơn vị đề-xi-mét vuông. (Nhập kết quả dưới dạng số thập phân).
\shortans{54,6}
\loigiai{
\textbf{Đáp số:} $54{,}6$\\
$8{,}4 \times 6{,}5 = 54{,}6\text{ (dm}^2\text{)}$.
}
\end{ex}

\begin{ex} %[5D4V24]
Trong hai năm trang trại bán được $21{,}56\text{ tấn}$ cá chẽm. Năm thứ hai bán nhiều hơn năm thứ nhất $2{,}7\text{ tấn}$. Hỏi năm thứ hai trang trại bán được bao nhiêu tấn cá chẽm? (Nhập kết quả dưới dạng số thập phân).
\shortans{12,13}
\loigiai{
\textbf{Đáp số:} $12{,}13$\\
Năm thứ hai bán được: $(21{,}56 + 2{,}7) : 2 = 12{,}13\text{ (tấn)}$.
}
\end{ex}

\begin{ex} %[5D4V24]
Khu vườn hình chữ nhật có diện tích $83{,}52\text{ m}^2$ và chiều rộng $8{,}7\text{ m}$. Tính chu vi khu vườn đó theo mét. (Nhập kết quả dưới dạng số thập phân).
\shortans{36,6}
\loigiai{
\textbf{Đáp số:} $36{,}6$\\
Chiều dài: $83{,}52 : 8{,}7 = 9{,}6\text{ m}$. Chu vi: $(9{,}6 + 8{,}7) \times 2 = 36{,}6\text{ m}$.
}
\end{ex}
